\newpage
\section*{Abstract}

This work presents a low-cost 3D scanner built only with an ESP32
microcontroller and a VL53L0X time-of-flight laser sensor. The sensor sits
inside a box that looks downward and slides at a fixed height over a grid
drawn on the table. With a single button, a marker sweeps each row at constant
speed and the user slides the box after it; each reading is located with an
affine function of time and averaged in its cell.
A single $4\times4$ homogeneous matrix turns each cell and its distance into a
point in space, combining a scaling, a reflection and a translation. The tilt
of the table is estimated by fitting a plane with least squares to the border
cells, and the resulting surface is drawn in a web interface with 3D
rotations, orthogonal projection and shading computed from cross and dot
products. The interface breaks down these calculations with the numbers of
each measurement. In simulation, with an empty tilted table, the plane fit
reduced the height error from \simMesaSin{} mm to \simMesaCon{} mm, and the
height of a 95 mm mug was measured as \simTazaMax{} mm. With the real sensor,
ten measurements at the same position varied with a standard deviation of
\repDesv{} mm. The sensor's cone of about $25^\circ$ smooths object edges,
which limits lateral resolution. The prototype's parts cost less than eighty
thousand Colombian pesos.

\noindent\textit{Keywords:} linear algebra, homogeneous matrix, least squares,
3D rotations, cross product, 3D scanner.
\newpage
